% 大阪大学基礎工学部知能システム学コース 特別研究概要スタイルファイル
\typeout{大阪大学基礎工学部知能システム学コース 特別研究概要スタイルファイル(2019/12/24) by 堀井隆斗}

\usepackage{amsmath,amssymb,bm}
\usepackage[dvipdfmx]{graphicx}
%\usepackage[dvips]{graphicx}
\usepackage{ascmac}
\usepackage{calc}
\usepackage{kvsetkeys}
\usepackage{kvoptions}
\usepackage{tikz}
\usepackage{progressbar}
\usepackage{algorithm}
\usepackage{algpseudocode}
\usepackage {subfigmat}
\usepackage {diagbox}
\flushbottom
\sloppy

% 余白

\setlength{\paperwidth}{210mm}
\setlength{\paperheight}{297mm}
\setlength{\voffset}{0mm}
\setlength{\textwidth}{\paperwidth}
\addtolength{\textwidth}{-30mm}
\setlength{\textheight}{\paperheight}
\addtolength{\textheight}{-40mm}
\setlength{\topmargin}{-1in}
\addtolength{\topmargin}{15mm}
\setlength{\headheight}{0mm}
\setlength{\headsep}{0mm}
\setlength{\footskip}{10mm}
\setlength{\oddsidemargin}{-1in}
\addtolength{\oddsidemargin}{15mm}
\setlength{\columnsep}{7mm}

% ヘッダ＆タイトル

\makeatletter
\def\headding#1{\gdef\@headding{#1}}
%\headding{研究進捗報告書}
\def\@maketitle{\newpage
\setlength{\topsep}{0mm}
 \null
 \begin{center}
 {\Large {\bf \@title} \par}
 \vspace{1mm}
 {\large 学籍番号：{\@date} \hspace{10pt} 吉川研究室 \hspace{10pt} {\@author} \par}
 \end{center}
 \vskip 1.5\baselineskip
 }
% 章間の間隔を調整する
\renewcommand{\section}{%
  \@startsection{section}{1}{\z@}%
  {0.4\Cvs}{0.1\Cvs}%
  {\large\bf}}
%  {\normalfont\large\headfont\raggedright}}

\renewcommand{\subsection}{%
  \@startsection{subsection}{1}{\z@}%
  {0.4\Cvs}{0.1\Cvs}%
  {\large\bf}}

\renewcommand{\subsubsection}{%
  \@startsection{subsubsection}{1}{\z@}%
  {0.4\Cvs}{0.1\Cvs}%
  {\normalsize}}
\makeatother

% 図のキャプションをFig.に，表のキャプションをTableにする

\def\fnum@figure{{\bf \figurename}~\thefigure}
\renewcommand{\figurename}{Fig.}

\def\fnum@table{{\bf \tablename~\thetable}}
\renewcommand{\tablename}{Table}

% 箇条書きの行間調節

\renewenvironment{itemize}{%
  \begin{list}{\theenumi.~}{%
  \usecounter{enumi}%
  \setlength{\itemsep}{-.5\baselineskip}%
  \setlength{\parsep}{.5\baselineskip}%
  }}{\end{list}}

% 本文中の数式を別行立て数式と同じように表示する

\everymath{\displaystyle}

\pagestyle{empty} % すべてのページ番号を消